% Luc Destombe --- licence CC BY-NC-SA 4.0
% https://creativecommons.org/licenses/by-nc-sa/4.0/deed.fr
% Fichier autonome : compiler avec pdflatex, deux fois (signets, nombre de pages).
\documentclass[11pt,a4paper]{article}
\makeatletter
\newif\iflycee@unecolonne
\lycee@unecolonnetrue
\newif\iflycee@palatino
\lycee@palatinotrue

\RequirePackage[utf8]{inputenc}
\RequirePackage[T1]{fontenc}
\RequirePackage[french]{babel}
\RequirePackage{etoolbox}

\newcommand{\ShorthandsOff}{\@ifundefined{shorthandoff}{}{\shorthandoff{;:!?}}}
\newcommand{\ShorthandsOn}{\@ifundefined{shorthandon}{}{\shorthandon{;:!?}}}
\ShorthandsOff
\AtBeginDocument{\ShorthandsOn}
\AtBeginEnvironment{tikzpicture}{\ShorthandsOff}

\RequirePackage{geometry}
\RequirePackage{fancyhdr}
\RequirePackage{lastpage}
\RequirePackage{multicol}
\RequirePackage{array}
\RequirePackage{enumitem}
\RequirePackage{xcolor}
\RequirePackage{needspace}

\RequirePackage{amsmath,amssymb}
\RequirePackage{mathpazo}
\DeclareMathSizes{10.95}{10.95}{8}{5.5}
\begingroup\catcode`\_=\active \catcode`\^=\active
\gdef_#1{\sb{\scriptscriptstyle #1}}
\gdef^#1{\sp{\scriptscriptstyle\lycee@moinsepais #1}}
\endgroup
\newcommand\lycee@moins{\mathbin{%
  \setbox\z@\hbox{$\m@th\scriptscriptstyle\lycee@moinsorig$}%
  \dimen@=.55\wd\z@ \advance\dimen@-.5pt
  \hbox to.75\wd\z@{\hss$\m@th\scriptscriptstyle\vcenter{\hbox{%
    \tikz\draw[line width=.5pt,line cap=round](0,0)--(\the\dimen@,0);}}$\hss}}}
\begingroup\lccode`\~=`\- \lowercase{\endgroup\let~}\lycee@moins
\newcommand\lycee@moinsepais{\mathcode`\-="8000 }
\AtBeginDocument{\mathchardef\lycee@moinsorig=\mathcode`\-\relax}
\AtBeginDocument{\catcode`\_=12 \mathcode`\_="8000
  \catcode`\^=12 \mathcode`\^="8000 }
\newcommand\lycee@exposants{%
  \fontdimen13\textfont2=.48\fontdimen6\textfont2
  \fontdimen14\textfont2=.48\fontdimen6\textfont2
  \fontdimen15\textfont2=.48\fontdimen6\textfont2
  \fontdimen13\scriptfont2=.48\fontdimen6\scriptfont2
  \fontdimen14\scriptfont2=.48\fontdimen6\scriptfont2
  \fontdimen15\scriptfont2=.48\fontdimen6\scriptfont2}
\every@math@size\expandafter{\the\every@math@size\lycee@exposants}
\newlength{\retraitcalculs}
\setlength{\retraitcalculs}{2em}

\RequirePackage{tikz}
\usetikzlibrary{arrows.meta,calc,babel,shapes.misc}
\RequirePackage[most]{tcolorbox}
\tcbuselibrary{skins,breakable}

\definecolor{coulDef}{HTML}{1F4E79}
\definecolor{coulProp}{HTML}{7030A0}
\definecolor{coulRem}{HTML}{C55A11}
\definecolor{coulEx}{HTML}{2E7D32}
\definecolor{coulExo}{HTML}{B03A2E}
\definecolor{coulPart}{HTML}{1F4E79}
\definecolor{coulMeth}{HTML}{1B6E4F}
\definecolor{coulCourbe}{HTML}{0B5ED7}
\definecolor{coulCourbeB}{HTML}{C00000}
\definecolor{coulCourbeC}{HTML}{2E7D32}
\definecolor{coulGrille}{HTML}{8C8C8C}
\definecolor{coulRep}{HTML}{1B6E4F}
\definecolor{coulBar}{HTML}{2E7D32}

\newcommand{\R}{\mathbb{R}}
\newcommand{\Rs}{\mathbb{R}^{*}}
\renewcommand{\le}{\leqslant}
\renewcommand{\ge}{\geqslant}
\newcommand{\vect}[1]{\overrightarrow{#1}}
\newcommand{\norme}[1]{\left\lVert #1 \right\rVert}
\newcommand{\intervalle}[2]{\left[#1\,;#2\right]}
\RequirePackage{cancel}
\renewcommand{\CancelColor}{\color{coulExo}}
\newcommand{\fois}[1]{\times{\color{coulExo}#1}}
\newcommand{\barre}[1]{\cancel{#1}}

\tikzset{grille/.style={coulGrille,line width=0.4pt}}
\newcommand{\quadrillage}[4]{\draw[grille,step=1] (#1,#2) grid (#3,#4);}
\tikzset{
  croix/.style={cross out,draw,line width=0.6pt,inner sep=0pt,minimum size=4pt},
  croixdroite/.style={croix,rotate=45,line width=0.8pt},
}
\newcommand{\pt}[3]{\path (#1) node[croix]{} node[#2,font=\small,inner sep=2.5pt]{$#3$};}

\newcommand{\lycee@repere}[4]{%
  \draw[grille,step=1] (#1,#2) grid (#3,#4);
  \draw[-{Stealth[length=2mm]},thick] (#1,0)--({#3+0.4},0) node[below left=-1pt]{$x$};
  \draw[-{Stealth[length=2mm]},thick] (0,#2)--(0,{#4+0.4}) node[below left=-1pt]{$y$};
  \node[below left,font=\scriptsize,inner sep=1.5pt] at (0,0) {$0$};}
\newcommand{\lycee@gradx}[1]{\foreach \k/\l in {#1}{%
  \draw (\k,0.07)--(\k,-0.07) node[below,font=\scriptsize,inner sep=1.5pt]{$\l$};}}
\newcommand{\lycee@grady}[1]{\foreach \k/\l in {#1}{%
  \draw (0.07,\k)--(-0.07,\k) node[left,font=\scriptsize,inner sep=1.5pt]{$\l$};}}
\newcommand{\lycee@intff}[2]{\left[#1\,;#2\right]}
\newcommand{\lycee@intfo}[2]{\left[#1\,;#2\right[}
\newcommand{\lycee@intof}[2]{\left]#1\,;#2\right]}
\newcommand{\lycee@intoo}[2]{\left]#1\,;#2\right[}
\newcommand{\lycee@ens}[1]{\left\{#1\right\}}
\AtBeginDocument{%
  \providecommand{\repere}{\lycee@repere}%
  \providecommand{\gradx}{\lycee@gradx}%
  \providecommand{\grady}{\lycee@grady}%
  \providecommand{\Cf}{\mathcal{C}\sb{\scriptscriptstyle f}}%
  \providecommand{\Cg}{\mathcal{C}\sb{\scriptscriptstyle g}}%
  \providecommand{\intff}{\lycee@intff}%
  \providecommand{\intfo}{\lycee@intfo}%
  \providecommand{\intof}{\lycee@intof}%
  \providecommand{\intoo}{\lycee@intoo}%
  \providecommand{\ens}{\lycee@ens}}

\newlist{questions}{enumerate}{3}
\setlist[questions,1]{label=\arabic*\textdegree),leftmargin=*,itemsep=2pt,topsep=3pt}
\setlist[questions,2]{label=\alph*),leftmargin=*,itemsep=1pt,topsep=2pt}
\setlist[questions,3]{label=\roman*),leftmargin=*,itemsep=1pt,topsep=1pt}
\setlist[itemize]{label=\textbullet,leftmargin=*,itemsep=2pt,topsep=3pt}

\newcommand{\besoinplace}[1]{\par\penalty\@M\begingroup
  \dimen@=#1\relax\dimen@ii=\pagegoal\advance\dimen@ii-\pagetotal
  \ifdim\dimen@>\dimen@ii\ifdim\dimen@ii>\z@\vfil\fi\break\fi\endgroup}

\newcommand{\lycee@pied}{\thepage/\pageref{LastPage}}
\newcommand{\entete}[2]{%
  \pagestyle{fancy}\fancyhf{}%
  \fancyhead[L]{\footnotesize\color{coulPart}\textbf{#1}}%
  \fancyhead[R]{\footnotesize\color{coulPart}\textbf{#2}}%
  \fancyfoot[C]{\footnotesize\lycee@pied}%
  \renewcommand{\headrulewidth}{0.4pt}%
  \renewcommand{\headrule}{\hbox to\headwidth{%
    \color{coulPart!40}\leaders\hrule height \headrulewidth\hfill}}}

\RequirePackage{ccicons}
\newcommand{\lycee@licence}{{\scriptsize\color{gray}%
  \copyright\ Luc Destombe --- \ccbyncsa\ CC BY-NC-SA 4.0}}
\newif\iflycee@licence \lycee@licencetrue
\newcommand{\sanslicence}{\lycee@licencefalse}
\AtBeginDocument{\iflycee@licence\AddToHookNext{shipout/foreground}{%
  \put(0,0){\rlap{\kern\dimexpr1in+\hoffset+\oddsidemargin+\textwidth\relax
    \llap{\raisebox{-\dimexpr1in+\voffset+\topmargin+\headheight+\headsep
      +\textheight+\footskip\relax}{\lycee@licence}}}}}\fi}

\newcommand{\formulesserrees}{%
  \setlength{\abovedisplayskip}{3pt}\setlength{\belowdisplayskip}{3pt}%
  \setlength{\abovedisplayshortskip}{2pt}\setlength{\belowdisplayshortskip}{3pt}}

\setlength{\parindent}{0pt}

\RequirePackage[implicit=false,hidelinks,pdfencoding=auto,psdextra,bookmarksopen,
  bookmarksopenlevel=1,pdfcreator={},pdfproducer={}]{hyperref}
\RequirePackage{bookmark}
\hypersetup{pdfauthor={Luc Destombe},pdfkeywords={CC BY-NC-SA 4.0}}
\newcounter{lycee@signet}
\newcommand{\signet}[3][]{\stepcounter{lycee@signet}%
  \edef\lycee@dest{\if\relax\detokenize{#1}\relax
    signet.\arabic{lycee@signet}\else#1\fi}%
  \hypertarget{\lycee@dest}{}\bookmark[level=#2,dest=\lycee@dest]{#3}}
\pdfstringdefDisableCommands{%
  \def\R{R}\def\vect#1{#1}\def\Cf{Cf}\def\Cg{Cg}\def\fois#1{×#1}%
  \def\barre#1{#1}\def\intervalle#1#2{[#1;#2]}\def\intff#1#2{[#1;#2]}%
  \def\textsuperscript#1{#1}\def\dfrac#1#2{#1/#2}\def\frac#1#2{#1/#2}%
  \def\vec#1{vecteur #1}}
\RequirePackage[official]{eurosym}

\tcbset{exercice corrige/.style={
  enhanced,breakable,before upper=\raggedright,
  colback=white,colframe=coulExo,
  boxrule=0.8pt,arc=2pt,
  left=6pt,right=6pt,top=8pt,bottom=5pt,
  before skip=9pt,after skip=4pt,
  coltitle=white,fonttitle=\bfseries\small,
  attach boxed title to top left={xshift=8pt,yshift=-\tcboxedtitleheight/2},
  boxed title style={colback=coulExo,arc=2pt,boxrule=0pt}}}
\newcounter{exo}
\newtcolorbox{exercice}[1][]{exercice corrige,
  phantom={\signet[exercice-\arabic{exo}]{2}{Exercice \arabic{exo}}},
  title={Exercice \theexo\ifx\relax#1\relax\else{} --- #1\fi}}
\newenvironment{exo}[1][\relax]{\refstepcounter{exo}\begin{exercice}[#1]}{\end{exercice}}

\geometry{top=1.8cm,bottom=1cm,left=1cm,right=1cm,headheight=14pt,footskip=0.6cm}
\renewcommand{\lycee@pied}{\thepage}

\renewcommand{\theexo}{\ifnum\value{exo}<10 0\fi\arabic{exo}}
\tcbset{exercice corrige/.style={
  enhanced,breakable,before upper=\raggedright,
  colback=white,colframe=coulExo,
  boxrule=0.9pt,arc=2pt,
  left=8pt,right=8pt,top=8pt,bottom=6pt,
  before skip=8pt,after skip=6pt,
  coltitle=white,fonttitle=\bfseries,
  attach boxed title to top left={xshift=8pt,yshift=-\tcboxedtitleheight/2},
  boxed title style={colback=coulExo,arc=2pt,boxrule=0pt}}}

\newcommand{\entetefiche}[2]{%
  \begin{center}
    \begin{tcolorbox}[enhanced,width=0.92\linewidth,colback=white,colframe=coulPart,
        boxrule=1.4pt,arc=3pt,halign=center,top=6pt,bottom=6pt]
      {\large\bfseries\color{coulPart} CORRIGÉ --- FICHE D'EXERCICES}\\[3pt]
      {\normalsize\bfseries #1}\\[2pt]
      {\normalsize #2}
    \end{tcolorbox}
  \end{center}}

\newcounter{partie}
\newcommand{\partie}[1]{%
  \stepcounter{partie}%
  \besoinplace{8\baselineskip}\vspace{4pt}%
  \begin{tcolorbox}[enhanced,colback=coulPart!8,colframe=coulPart,boxrule=0pt,
      leftrule=3pt,arc=0pt,left=6pt,right=4pt,top=3pt,bottom=3pt,
      before skip=6pt,after skip=2pt]
    \bfseries\color{coulPart}\Roman{partie} --- #1\signet{1}{\Roman{partie} --- #1}
  \end{tcolorbox}}

\setlist[questions,1]{label=\arabic*\textdegree),leftmargin=*,itemsep=3pt,topsep=3pt}
\setlist[questions,2]{label=\alph*),leftmargin=*,itemsep=2pt,topsep=2pt}
\setlist[questions,3]{label=\roman*),leftmargin=*,itemsep=1pt,topsep=1pt}
\newlist{expr}{enumerate}{1}
\setlist[expr]{label={\Alph*\ $=$},leftmargin=*,itemsep=2pt,topsep=3pt}

\newcommand{\res}[1]{{\color{coulRep}#1}}
\newcommand{\rem}[1]{\par\smallskip{\small\itshape\color{black!60}$\rightarrow$ #1}}
\newenvironment{calculs}{\par\smallskip\noindent$\displaystyle\begin{aligned}[t]}%
  {\end{aligned}$\par\smallskip}
\newcommand{\justif}[1]{\quad\text{\small\itshape\color{black!60}#1}}

\setlength{\parskip}{2pt}

\makeatother
\geometry{top=1.8cm,bottom=1.6cm,left=1.8cm,right=1.8cm,headheight=14pt,footskip=30pt}
\entete{Chapitre 1 --- Calcul littéral --- corrigé}{Première spécialité}

\usepackage{tkz-tab}

\begin{document}

\entetefiche{Chapitre 1 --- Calcul littéral}{Première spécialité mathématiques}

\partie{Réduire, développer, identités remarquables}

\begin{exo}[Réduire et distribuer]
\begin{multicols}{2}
\begin{expr}
  \item $3x+8+6x=\res{9x+8}$
  \item $4x^{2}-9x^{2}+2x^{2}=\res{-3x^{2}}$
  \item $6(3x-2)=\res{18x-12}$
  \item $2x(5x-4)=\res{10x^{2}-8x}$
  \item $9x-9-4x=\res{5x-9}$
  \item $-5x+6x+21=\res{x+21}$
  \item $-(4-3x)=\res{3x-4}$
  \item $-2x(1-4x)=\res{8x^{2}-2x}$
\end{expr}
\end{multicols}
\rem{En G et H, le « $-$ » devant la parenthèse change le signe de \emph{chacun} de ses
termes.}
\end{exo}

\begin{exo}[Double distributivité et identités remarquables]
\begin{multicols}{2}
\begin{expr}
  \item $8x^{2}+2x+12x+3=\res{8x^{2}+14x+3}$
  \item $-6x^{2}+15x+8x-20=\res{-6x^{2}+23x-20}$
  \item $(5x)^{2}+2\times 5x\times 2+2^{2}=\res{25x^{2}+20x+4}$
  \item $4^{2}-2\times 4\times 3x+(3x)^{2}=\res{9x^{2}-24x+16}$
  \item $(3x)^{2}-5^{2}=\res{9x^{2}-25}$
  \item $x^{2}-2\times x\times\frac{1}{3}+\frac{1}{9}=\res{x^{2}-\frac{2}{3}x+\frac{1}{9}}$
\end{expr}
\end{multicols}
\begin{expr}[start=7]
  \item $3(x+2)(1-x)$
        \begin{calculs}
        G &= 3\left(x-x^{2}+2-2x\right) &&\justif{on développe d'abord les deux parenthèses}\\
        G &= 3\left(-x^{2}-x+2\right)\\
        G &= \res{-3x^{2}-3x+6}
        \end{calculs}
\end{expr}
\rem{De C à F, l'identité remarquable donne le résultat sans tout distribuer.}
\end{exo}

\begin{exo}[Développer, réduire, ordonner]
\begin{expr}
  \item $(4x+1)(2x+3)-(5x-2)(x+4)$
        \begin{calculs}
        A &= \left(8x^{2}+14x+3\right)-\left(5x^{2}+18x-8\right)\\
        A &= 8x^{2}+14x+3-5x^{2}-18x+8\\
        A &= \res{3x^{2}-4x+11}
        \end{calculs}
  \item $(2x-5)^{2}-(3x+1)(x-2)$
        \begin{calculs}
        B &= \left(4x^{2}-20x+25\right)-\left(3x^{2}-5x-2\right)\\
        B &= \res{x^{2}-15x+27}
        \end{calculs}
  \item $\left(x^{2}-x+2\right)(3x+1)-(x+1)(x-1)^{2}$
        \begin{calculs}
        \left(x^{2}-x+2\right)(3x+1) &= 3x^{3}+x^{2}-3x^{2}-x+6x+2\\
        \left(x^{2}-x+2\right)(3x+1) &= 3x^{3}-2x^{2}+5x+2\\
        (x+1)(x-1)^{2} &= \left(x^{2}-1\right)(x-1) &&\justif{$(x+1)(x-1)=x^{2}-1$}\\
        (x+1)(x-1)^{2} &= x^{3}-x^{2}-x+1\\
        C &= \res{2x^{3}-x^{2}+6x+1}
        \end{calculs}
  \item $(x+2)\left(x^{2}-2x+4\right)$
        \begin{calculs}
        D &= x^{3}-2x^{2}+4x+2x^{2}-4x+8\\
        D &= \res{x^{3}+8}
        \end{calculs}
  \item $3x(2-x)-(x-4)(2x+3)$
        \begin{calculs}
        E &= \left(6x-3x^{2}\right)-\left(2x^{2}-5x-12\right)\\
        E &= \res{-5x^{2}+11x+12}
        \end{calculs}
\end{expr}
\rem{Le piège commun : le « $-$ » devant le second produit. On garde les parenthèses
jusqu'au bout du développement.}
\end{exo}

\partie{Substituer une valeur}

\begin{exo}[Calculer une image]
\begin{questions}
  \item Avec $P(x)=3x^{2}-4x-2$ :
        \begin{questions}[label=\alph*)]
          \item $P(1)=3-4-2=\res{-3}$ \quad et \quad $P(-1)=3+4-2=\res{5}$
          \item $P\!\left(\frac{1}{2}\right)=\frac{3}{4}-2-2=\res{-\frac{13}{4}}$ \quad et
                \quad $P\!\left(-\frac{2}{3}\right)=\frac{4}{3}+\frac{8}{3}-2=\res{2}$
        \end{questions}
  \item \begin{questions}[label=\alph*)]
          \item Pour $x=0$ : $A=2\times(-5)$, soit $A=\res{-10}$, et $B=\res{-10}$.
          \item Pour $x=1$ : $A=3\times(-4)$, soit $A=\res{-12}$, et $B=1+3-10$, soit
                $B=\res{-6}$.
          \item \res{Non} : pour $x=1$, $A$ et $B$ ont deux valeurs différentes, donc ils
                ne sont pas égaux pour tout réel $x$. Une seule valeur commune ne prouve
                rien.
        \end{questions}
\end{questions}
\rem{En développant, $A=x^{2}-3x-10$, qui diffère de $B$ par le terme en $x$.}
\end{exo}

\begin{exo}[Substituer une racine carrée]
\begin{questions}[label=\alph*)]
  \item $P\!\left(\sqrt{5}\right)=5-5=\res{0}$
  \item $P\!\left(\sqrt{2}\right)=2\times 2+3\sqrt{2}-1=\res{3+3\sqrt{2}}$
  \item \begin{calculs}
        P\!\left(\sqrt{3}\right) &= \left(3\sqrt{3}+1\right)\left(\sqrt{3}-2\right)\\
        P\!\left(\sqrt{3}\right) &= 9-6\sqrt{3}+\sqrt{3}-2 &&\justif{$3\sqrt{3}\times\sqrt{3}=3\times 3$}\\
        P\!\left(\sqrt{3}\right) &= \res{7-5\sqrt{3}}
        \end{calculs}
  \item \begin{calculs}
        P\!\left(\sqrt{5}\right) &= \left(2\sqrt{5}-1\right)^{2}+4\sqrt{5}\\
        P\!\left(\sqrt{5}\right) &= 20-4\sqrt{5}+1+4\sqrt{5} &&\justif{$\left(2\sqrt{5}\right)^{2}=4\times 5$}\\
        P\!\left(\sqrt{5}\right) &= \res{21}
        \end{calculs}
  \item \begin{calculs}
        P\!\left(1+\sqrt{2}\right) &= \left(1+\sqrt{2}\right)^{2}-2\left(1+\sqrt{2}\right)+3\\
        P\!\left(1+\sqrt{2}\right) &= 3+2\sqrt{2}-2-2\sqrt{2}+3 &&\justif{$\left(1+\sqrt{2}\right)^{2}=1+2\sqrt{2}+2$}\\
        P\!\left(1+\sqrt{2}\right) &= \res{4}
        \end{calculs}
\end{questions}
\rem{Le réflexe : $\left(\sqrt{a}\right)^{2}=a$. Les produits avec des racines se
développent comme les autres.}
\end{exo}

\begin{exo}[Racines carrées dans un quotient]
\begin{questions}
  \item $f(x)=\dfrac{x+1}{x-3}$.
        \begin{questions}[label=\alph*)]
          \item $x-3=0$ pour $x=3$ : la valeur interdite est \res{3}.
          \item $f(5)=\dfrac{6}{2}=\res{3}$ \quad et \quad
                $f\!\left(\frac{1}{2}\right)=\dfrac{\frac{3}{2}}{-\frac{5}{2}}=\res{-\dfrac{3}{5}}$
          \item On multiplie par la quantité conjuguée $\sqrt{5}+3$ :
                \begin{calculs}
                f\!\left(\sqrt{5}\right) &= \dfrac{\left(\sqrt{5}+1\right)\left(\sqrt{5}+3\right)}
                  {\left(\sqrt{5}-3\right)\left(\sqrt{5}+3\right)}\\[2pt]
                f\!\left(\sqrt{5}\right) &= \dfrac{5+3\sqrt{5}+\sqrt{5}+3}{5-9}\\[2pt]
                f\!\left(\sqrt{5}\right) &= \dfrac{8+4\sqrt{5}}{-4}\\[2pt]
                f\!\left(\sqrt{5}\right) &= \res{-2-\sqrt{5}}
                \end{calculs}
        \end{questions}
  \item Quantité conjuguée $\sqrt{3}-1$ :
        \begin{calculs}
        g\!\left(\sqrt{3}\right) &= \dfrac{\left(\sqrt{3}-2\right)\left(\sqrt{3}-1\right)}
          {\left(\sqrt{3}+1\right)\left(\sqrt{3}-1\right)}\\[2pt]
        g\!\left(\sqrt{3}\right) &= \dfrac{3-\sqrt{3}-2\sqrt{3}+2}{3-1}\\[2pt]
        g\!\left(\sqrt{3}\right) &= \res{\dfrac{5-3\sqrt{3}}{2}}
        \end{calculs}
  \item Le dénominateur vaut $\left(2+\sqrt{3}\right)-2=\sqrt{3}$ : on multiplie par
        $\sqrt{3}$.
        \begin{calculs}
        h\!\left(2+\sqrt{3}\right) &= \dfrac{5+\sqrt{3}}{\sqrt{3}}\\[2pt]
        h\!\left(2+\sqrt{3}\right) &= \dfrac{\left(5+\sqrt{3}\right)\sqrt{3}}{\sqrt{3}\times\sqrt{3}}\\[2pt]
        h\!\left(2+\sqrt{3}\right) &= \res{\dfrac{5\sqrt{3}+3}{3}}
        \end{calculs}
\end{questions}
\rem{Deux cas : un dénominateur qui est une \emph{somme} avec une racine se multiplie par
sa quantité conjuguée ; une racine \emph{seule} se multiplie par elle-même.}
\end{exo}

\partie{Factoriser une expression}

\begin{exo}[Facteur commun]
\begin{multicols}{2}
\begin{expr}
  \item $5x+20=\res{5(x+4)}$
  \item $7x^{2}+3x=\res{x(7x+3)}$
  \item $18x^{2}-12x=\res{6x(3x-2)}$
  \item $x^{2}+x=\res{x(x+1)}$
  \item $4x^{2}-x=\res{x(4x-1)}$
  \item $x(2+x+5)=\res{x(x+7)}$
\end{expr}
\end{multicols}
\rem{En D et E, le terme $x$ s'écrit $x\times 1$ : c'est le $1$ qu'on oublie. En F,
le premier terme $2x$ est $x\times 2$.}
\end{exo}

\begin{exo}[Identités remarquables et facteur commun]
\begin{expr}
  \item $(3x)^{2}+2\times 3x\times 5+5^{2}=\res{(3x+5)^{2}}$
  \item $(7x)^{2}-6^{2}=\res{(7x-6)(7x+6)}$
  \item $(3x+1)\bigl[(2x-5)+(x+7)\bigr]=\res{(3x+1)(3x+2)}$
  \item $(4x-1)\bigl[(x+6)-(3x-2)\bigr]=\res{(4x-1)(-2x+8)}$
        \quad (ou $2(4x-1)(4-x)$)
  \item $(3x+2)^{2}-(x-4)^{2}$
        \begin{calculs}
        E &= \bigl[(3x+2)-(x-4)\bigr]\bigl[(3x+2)+(x-4)\bigr] &&\justif{$a^{2}-b^{2}$}\\
        E &= (2x+6)(4x-2)\\
        E &= \res{4(x+3)(2x-1)} &&\justif{$2$ en facteur dans chaque parenthèse}
        \end{calculs}
  \item $36-(x+2)^{2}$
        \begin{calculs}
        F &= \bigl[6-(x+2)\bigr]\bigl[6+(x+2)\bigr]\\
        F &= \res{(4-x)(x+8)}
        \end{calculs}
\end{expr}
\rem{En D et E, le « $-$ » devant la seconde parenthèse change ses deux signes :
$-(3x-2)=-3x+2$.}
\end{exo}

\begin{exo}[Faire apparaître un facteur commun]
\begin{expr}
  \item On factorise d'abord $9x^{2}-16=(3x-4)(3x+4)$.
        \begin{calculs}
        A &= (3x-4)(3x+4)-(3x+4)(x-2)\\
        A &= (3x+4)\bigl[(3x-4)-(x-2)\bigr]\\
        A &= (3x+4)(2x-2)\\
        A &= \res{2(3x+4)(x-1)}
        \end{calculs}
  \item On factorise d'abord $x^{2}-25=(x-5)(x+5)$.
        \begin{calculs}
        B &= (x+5)\bigl[(x-5)-(2x-3)\bigr]\\
        B &= (x+5)(-x-2)\\
        B &= \res{-(x+5)(x+2)}
        \end{calculs}
  \item On écrit $x^{2}-9=(x-3)(x+3)$.
        \begin{calculs}
        C &= (x+3)\bigl[(4x-1)-(x-3)\bigr]\\
        C &= \res{(x+3)(3x+2)}
        \end{calculs}
  \item $x^{2}-\left(\sqrt{7}\right)^{2}=\res{\left(x-\sqrt{7}\right)\left(x+\sqrt{7}\right)}$
  \item $x^{2}+2\times x\times\sqrt{5}+\left(\sqrt{5}\right)^{2}=\res{\left(x+\sqrt{5}\right)^{2}}$
\end{expr}
\rem{De A à C, la différence de carrés cache le facteur commun. En D et E, $\sqrt{7}$ et
$\sqrt{5}$ jouent le rôle de $b$ dans l'identité remarquable.}
\end{exo}

\partie{Équations}

\begin{exo}[Équations du premier degré]
\begin{questions}[label=\alph*)]
  \item $4x-3x=2+9$, donc $x=11$ : $S=\res{\{11\}}$.
  \item $5x=25$, donc $x=5$ : $S=\res{\{5\}}$.
  \item $3x-3x=-1-4$ donne $0=-5$, ce qui est faux : $S=\res{\varnothing}$.
  \item \begin{calculs}
        3(x-2)=2(-x+7) &\iff 3x-6=-2x+14\\
                       &\iff 5x=20\\
                       &\iff x=4
        \end{calculs}
        $S=\res{\{4\}}$.
  \item \begin{calculs}
        x-(2-5x)=4x+9 &\iff 6x-2=4x+9\\
                      &\iff 2x=11
        \end{calculs}
        $S=\res{\left\{\frac{11}{2}\right\}}$.
  \item $4x-4+3x+15=0$, soit $7x+11=0$ : $S=\res{\left\{-\frac{11}{7}\right\}}$.
\end{questions}
\rem{En c), les termes en $x$ disparaissent et il reste une égalité fausse : aucune
solution. Ne pas écrire $S=\{0\}$.}
\end{exo}

\begin{exo}[Produit nul et carrés]
\begin{multicols}{2}
\begin{questions}[label=\alph*)]
  \item $x-6=0$ ou $4+x=0$ : $S=\res{\{-4\,;6\}}$
  \item $5x=0$ ou $3x-7=0$ : $S=\res{\left\{0\,;\frac{7}{3}\right\}}$
  \item $2x+9=0$ : $S=\res{\left\{-\frac{9}{2}\right\}}$
  \item $13>0$ : $S=\res{\left\{-\sqrt{13}\,;\sqrt{13}\right\}}$
  \item $x^{2}=7$ : $S=\res{\left\{-\sqrt{7}\,;\sqrt{7}\right\}}$
  \item $3x+1=5$ ou $3x+1=-5$ : $S=\res{\left\{-2\,;\frac{4}{3}\right\}}$
\end{questions}
\end{multicols}
\rem{En d) et e), on garde les valeurs exactes. En f), le carré vaut $25>0$ : deux cas.}
\end{exo}

\begin{exo}[Mise en équation]
\begin{questions}
  \item \textbf{La boulangerie.}
        \emph{Inconnue :} soit $x$ le nombre de croissants. Il y a $2x$ pains au chocolat
        et $x-10$ chaussons.
        \begin{calculs}
        x+2x+(x-10)=150 &\iff 4x-10=150\\
                        &\iff 4x=160\\
                        &\iff x=40
        \end{calculs}
        \emph{Vérification :} $40+80+30=150$.
        \emph{Conclusion :} la boulangerie a vendu \res{40 croissants, 80 pains au chocolat
        et 30 chaussons}.
  \item \textbf{Du carré au rectangle.} Le rectangle mesure $x+6$ sur $x-4$ ; son aire
        égale celle du carré, $x^{2}$.
        \begin{calculs}
        (x+6)(x-4)=x^{2} &\iff x^{2}+2x-24=x^{2}\\
                         &\iff 2x=24\\
                         &\iff x=12
        \end{calculs}
        \emph{Vérification :} le rectangle mesure $18$ cm sur $8$ cm, d'aire $144\text{ cm}^{2}$,
        comme le carré de côté $12$ cm.
        \emph{Conclusion :} \res{$x=12$ cm}, et le rectangle mesure \res{$18$ cm sur $8$ cm}.
\end{questions}
\rem{En 2\textdegree), les termes en $x^{2}$ s'éliminent : l'équation est du premier
degré.}
\end{exo}

\partie{Fractions rationnelles}

\begin{exo}[Les quatre opérations]
Dans les quatre cas, la valeur interdite est \res{$0$}.
\begin{multicols}{2}
\begin{expr}
  \item $\dfrac{3x}{4x}+\dfrac{20}{4x}=\res{\dfrac{3x+20}{4x}}$
  \item $\dfrac{3x}{4x}-\dfrac{20}{4x}=\res{\dfrac{3x-20}{4x}}$
  \item $\res{\dfrac{15}{4x}}$
  \item $\dfrac{3}{4}\times\dfrac{x}{5}=\res{\dfrac{3x}{20}}$
\end{expr}
\end{multicols}
\rem{Un dénominateur commun pour l'addition et la soustraction ; pas pour le produit.
Diviser par $\frac{5}{x}$, c'est multiplier par $\frac{x}{5}$.}
\end{exo}

\begin{exo}[Valeurs interdites]
\begin{expr}
  \item Valeur interdite \res{$4$} ; pas de simplification.
  \item $x(x-2)=0$ pour $x=0$ et pour $x=2$ : valeurs interdites \res{$0$ et $2$}.
  \item $2x+5=0$ pour $x=-\frac{5}{2}$ : valeur interdite \res{$-\frac{5}{2}$}.
  \item Valeur interdite \res{$4$}. Pour $x\ne 4$ :
        $\dfrac{(x-4)(x+4)}{x-4}=\res{x+4}$.
  \item $x^{2}+4\ge 4$, donc le dénominateur ne s'annule jamais : \res{aucune} valeur
        interdite.
  \item Valeur interdite \res{$7$}.
\end{expr}
\rem{En D, $x+4$ a une valeur pour $x=4$, mais pas l'expression de départ : simplifier ne
supprime pas une valeur interdite.}
\end{exo}

\begin{exo}[Réduire au même dénominateur]
\begin{expr}
  \item $x\ne-3$ : $\dfrac{2(x+3)+7}{x+3}=\res{\dfrac{2x+13}{x+3}}$
  \item $x\ne-1$ : $\dfrac{4(x+1)+x-2}{x+1}=\res{\dfrac{5x+2}{x+1}}$
  \item $x\ne 1$ et $x\ne-3$ :
        \begin{calculs}
        C &= \dfrac{x(x+3)-2x(x-1)}{(x-1)(x+3)}\\[2pt]
        C &= \dfrac{x^{2}+3x-2x^{2}+2x}{(x-1)(x+3)}\\[2pt]
        C &= \res{\dfrac{-x^{2}+5x}{(x-1)(x+3)}}
        \end{calculs}
  \item $x\ne-1$ et $x\ne\frac{3}{2}$ :
        \begin{calculs}
        D &= \dfrac{4(2x-3)-3(x+1)}{(x+1)(2x-3)}\\[2pt]
        D &= \res{\dfrac{5x-15}{(x+1)(2x-3)}}
        \end{calculs}
  \item $x\ne 2$ :
        \begin{calculs}
        E &= \dfrac{3x(x-2)-(2x+1)}{x-2}\\[2pt]
        E &= \res{\dfrac{3x^{2}-8x-1}{x-2}}
        \end{calculs}
  \item $3x-12=3(x-4)$ : le dénominateur commun est $3(x-4)$, et $x\ne 4$.
        \begin{calculs}
        F &= \dfrac{3\times 2x-(x+3)}{3(x-4)}\\[2pt]
        F &= \res{\dfrac{5x-3}{3(x-4)}}
        \end{calculs}
\end{expr}
\rem{En F, prendre $(x-4)(3x-12)$ comme dénominateur commun alourdit tout : repérer
$3x-12=3(x-4)$ fait gagner du temps.}
\end{exo}

\begin{exo}[Simplifier, multiplier, diviser]
\begin{expr}
  \item $x\ne-1$ : $\dfrac{4(x+1)\times(x+1)}{4(x+1)\times 2}=\res{\dfrac{x+1}{2}}$
  \item $x\ne 5$ : $\dfrac{(5-x)(5+x)}{x-5}=\dfrac{-(x-5)(x+5)}{x-5}$, donc
        \res{B=-(x+5)}.
  \item $x\ne-2$ et $x\ne 2$ :
        \begin{calculs}
        C &= \dfrac{(x+2)^{2}-(x-2)^{2}}{(x-2)(x+2)}\\[2pt]
        C &= \dfrac{x^{2}+4x+4-\left(x^{2}-4x+4\right)}{(x-2)(x+2)}\\[2pt]
        C &= \res{\dfrac{8x}{(x-2)(x+2)}}
        \end{calculs}
  \item $x\ne-4$ et $x\ne-1$ :
        \begin{calculs}
        D &= \dfrac{(x-1)(x+1)(x-4)(x+4)}{(x+4)(x+1)}\\[2pt]
        D &= \res{(x-1)(x-4)}
        \end{calculs}
  \item $x^{2}+3x=x(x+3)$ ; on divise par $\frac{x+2}{x+3}$, qui doit exister et être non
        nul : $x\ne 0$, $x\ne-3$ et $x\ne-2$.
        \begin{calculs}
        E &= \dfrac{(x-2)(x+2)}{x(x+3)}\times\dfrac{x+3}{x+2}\\[2pt]
        E &= \res{\dfrac{x-2}{x}}
        \end{calculs}
  \item $x^{2}-9=(x-3)(x+3)$ : valeurs interdites \res{$-3$ et $3$} ; le numérateur n'a
        pas de facteur commun avec le dénominateur, pas de simplification.
\end{expr}
\rem{En B, $25-x^{2}$ et $x-5$ sont « à l'envers » : $5-x=-(x-5)$. En E, le diviseur
ajoute la valeur interdite $-2$.}
\end{exo}

\partie{Équations quotient}

\begin{exo}[Quotient nul]
\emph{Un quotient est nul si et seulement si son numérateur est nul, son dénominateur ne
l'étant pas.}
\begin{questions}[label=\alph*)]
  \item Valeur interdite $-2$. $4x-12=0$ pour $x=3$, qui n'est pas interdit :
        $S=\res{\{3\}}$.
  \item Valeur interdite $0$. $3x+6=0$ pour $x=-2$ : $S=\res{\{-2\}}$.
  \item Valeur interdite $-1$. Le numérateur s'annule pour $x=3$ et pour $x=-\frac{5}{2}$ :
        $S=\res{\left\{-\frac{5}{2}\,;3\right\}}$.
  \item $x^{2}+3\ge 3$ : aucune valeur interdite. $7-2x=0$ pour $x=\frac{7}{2}$ :
        $S=\res{\left\{\frac{7}{2}\right\}}$.
\end{questions}
\end{exo}

\begin{exo}[Se ramener à un quotient nul]
\begin{questions}[label=\alph*)]
  \item Valeur interdite $-4$.
        \begin{calculs}
        \dfrac{3x+2}{x+4}=2 &\iff \dfrac{3x+2-2(x+4)}{x+4}=0\\[2pt]
                            &\iff \dfrac{x-6}{x+4}=0
        \end{calculs}
        $S=\res{\{6\}}$.
  \item Valeur interdite $-2$.
        \begin{calculs}
        \dfrac{-3x}{x+2}=1 &\iff \dfrac{-3x-(x+2)}{x+2}=0\\[2pt]
                           &\iff \dfrac{-4x-2}{x+2}=0
        \end{calculs}
        $S=\res{\left\{-\frac{1}{2}\right\}}$.
  \item Valeur interdite $0$. Dénominateur commun $3x$ :
        \begin{calculs}
        \dfrac{1}{x}+\dfrac{1}{3}=1 &\iff \dfrac{3+x-3x}{3x}=0\\[2pt]
                                    &\iff \dfrac{3-2x}{3x}=0
        \end{calculs}
        $S=\res{\left\{\frac{3}{2}\right\}}$.
  \item Valeur interdite $6$.
        \begin{calculs}
        \dfrac{5}{x-6}+1=0 &\iff \dfrac{5+x-6}{x-6}=0\\[2pt]
                           &\iff \dfrac{x-1}{x-6}=0
        \end{calculs}
        $S=\res{\{1\}}$.
  \item Valeur interdite $-2$.
        \begin{calculs}
        3-\dfrac{8-x}{x+2}=0 &\iff \dfrac{3(x+2)-(8-x)}{x+2}=0\\[2pt]
                             &\iff \dfrac{4x-2}{x+2}=0
        \end{calculs}
        $S=\res{\left\{\frac{1}{2}\right\}}$.
  \item Valeur interdite $4$. $x^{2}-16=0$ pour $x=4$ ou $x=-4$, mais $4$ est
        \emph{interdit} : $S=\res{\{-4\}}$.
\end{questions}
\rem{En f), une solution du numérateur est une valeur interdite : on l'écarte.}
\end{exo}

\begin{exo}[Égalité de deux quotients]
\emph{On ramène tout dans un membre et on réduit au même dénominateur (méthode du
cours).}
\begin{questions}[label=\alph*)]
  \item Valeurs interdites $-2$ et $0$.
        \begin{calculs}
        \dfrac{5}{x+2}=\dfrac{3}{x} &\iff \dfrac{5x-3(x+2)}{x(x+2)}=0\\[2pt]
                                    &\iff \dfrac{2x-6}{x(x+2)}=0
        \end{calculs}
        $S=\res{\{3\}}$.
  \item Valeurs interdites $0$ et $-1$. Numérateur :
        \begin{calculs}
        (2x+1)(x+1)-x(2x+5) &= 2x^{2}+3x+1-2x^{2}-5x\\
        (2x+1)(x+1)-x(2x+5) &= -2x+1
        \end{calculs}
        $S=\res{\left\{\frac{1}{2}\right\}}$.
  \item Valeurs interdites $-1$ et $2$. Numérateur :
        \begin{calculs}
        (3x-2)(x-2)-(3x+1)(x+1) &= \left(3x^{2}-8x+4\right)-\left(3x^{2}+4x+1\right)\\
        (3x-2)(x-2)-(3x+1)(x+1) &= -12x+3
        \end{calculs}
        $S=\res{\left\{\frac{1}{4}\right\}}$.
  \item Valeur interdite $2$.
        \begin{calculs}
        \dfrac{x+4}{x-2}-3=\dfrac{8-x}{x-2}
          &\iff \dfrac{x+4-3(x-2)-(8-x)}{x-2}=0\\[2pt]
          &\iff \dfrac{-x+2}{x-2}=0
        \end{calculs}
        Le numérateur s'annule pour $x=2$, qui est \emph{interdit} : $S=\res{\varnothing}$.
  \item Valeur interdite $1$. Même dénominateur :
        $\dfrac{x+5-3}{x-1}=0$, soit $\dfrac{x+2}{x-1}=0$ : $S=\res{\{-2\}}$.
  \item Valeurs interdites $1$ et $3$. On réduit d'abord le membre de gauche :
        \begin{calculs}
        1-\dfrac{x+1}{x-1} &= \dfrac{x-1-(x+1)}{x-1}\\[2pt]
        1-\dfrac{x+1}{x-1} &= \dfrac{-2}{x-1}
        \end{calculs}
        puis
        \begin{calculs}
        \dfrac{-2}{x-1}=\dfrac{4}{3-x}
          &\iff \dfrac{-2(3-x)-4(x-1)}{(x-1)(3-x)}=0\\[2pt]
          &\iff \dfrac{-2x-2}{(x-1)(3-x)}=0
        \end{calculs}
        $S=\res{\{-1\}}$.
\end{questions}
\rem{En d), la seule valeur trouvée est interdite : pas de solution. En e), les
dénominateurs sont égaux : on compare directement les numérateurs.}
\end{exo}

\partie{Inéquations et tableaux de signes}

\begin{exo}[Inéquations du premier degré]
\begin{questions}[label=\alph*)]
  \item $3x>-12$, donc $x>-4$ : $S=\res{\left]-4\,;+\infty\right[}$.
  \item $5x<10$, donc $x<2$ : $S=\res{\left]-\infty\,;2\right[}$.
  \item $3x-3\ge 7-2x$, soit $5x\ge 10$ : $S=\res{\left[2\,;+\infty\right[}$.
  \item On multiplie par $3>0$ (le sens ne change pas) : $2-5x>0$, soit $-5x>-2$, puis
        on divise par $-5<0$ (le sens change) : $x<\frac{2}{5}$.
        $S=\res{\left]-\infty\,;\frac{2}{5}\right[}$.
  \item $3x+6<3x+4$ donne $6<4$, ce qui est faux : $S=\res{\varnothing}$.
  \item On multiplie par $4>0$ :
        \begin{calculs}
        \dfrac{x-3}{2}-\dfrac{1-x}{4}>0 &\iff 2(x-3)-(1-x)>0\\
                                        &\iff 3x-7>0
        \end{calculs}
        $S=\res{\left]\frac{7}{3}\,;+\infty\right[}$.
\end{questions}
\rem{Seuls la multiplication et la division par un nombre \emph{négatif} changent le
sens. En e), l'inégalité obtenue est fausse : aucune solution.}
\end{exo}

\begin{exo}[Tableaux de signes]
\begin{expr}
  \item $(x-5)(x+1)$
\begin{center}
\begin{tikzpicture}[scale=0.85]
  \tkzTabInit[lgt=2.6,espcl=2]{$x$/1,$x+1$/1,$x-5$/1,$(x-5)(x+1)$/1.2}{$-\infty$,$-1$,$5$,$+\infty$}
  \tkzTabLine{,-,z,+,t,+,}
  \tkzTabLine{,-,t,-,z,+,}
  \tkzTabLine{,+,z,-,z,+,}
\end{tikzpicture}
\end{center}
  \item $(3-x)(2x+1)$
\begin{center}
\begin{tikzpicture}[scale=0.85]
  \tkzTabInit[lgt=2.6,espcl=2]{$x$/1,$2x+1$/1,$3-x$/1,$(3-x)(2x+1)$/1.2}{$-\infty$,$-\frac{1}{2}$,$3$,$+\infty$}
  \tkzTabLine{,-,z,+,t,+,}
  \tkzTabLine{,+,t,+,z,-,}
  \tkzTabLine{,-,z,+,z,-,}
\end{tikzpicture}
\end{center}
  \item $x^{2}-16=(x-4)(x+4)$
\begin{center}
\begin{tikzpicture}[scale=0.85]
  \tkzTabInit[lgt=2.6,espcl=2]{$x$/1,$x+4$/1,$x-4$/1,$x^{2}-16$/1.2}{$-\infty$,$-4$,$4$,$+\infty$}
  \tkzTabLine{,-,z,+,t,+,}
  \tkzTabLine{,-,t,-,z,+,}
  \tkzTabLine{,+,z,-,z,+,}
\end{tikzpicture}
\end{center}
  \item $2x(x-3)(4x+1)$
\begin{center}
\begin{tikzpicture}[scale=0.85]
  \tkzTabInit[lgt=3,espcl=1.9]{$x$/1,$4x+1$/1,$2x$/1,$x-3$/1,$2x(x-3)(4x+1)$/1.2}{$-\infty$,$-\frac{1}{4}$,$0$,$3$,$+\infty$}
  \tkzTabLine{,-,z,+,t,+,t,+,}
  \tkzTabLine{,-,t,-,z,+,t,+,}
  \tkzTabLine{,-,t,-,t,-,z,+,}
  \tkzTabLine{,-,z,+,z,-,z,+,}
\end{tikzpicture}
\end{center}
  \item $\dfrac{x+4}{1-x}$ (valeur interdite $1$)
\begin{center}
\begin{tikzpicture}[scale=0.85]
  \tkzTabInit[lgt=2.6,espcl=2]{$x$/1,$x+4$/1,$1-x$/1,$\frac{x+4}{1-x}$/1.4}{$-\infty$,$-4$,$1$,$+\infty$}
  \tkzTabLine{,-,z,+,t,+,}
  \tkzTabLine{,+,t,+,z,-,}
  \tkzTabLine{,-,z,+,d,-,}
\end{tikzpicture}
\end{center}
  \item $\dfrac{x(x-2)}{2x+3}$ (valeur interdite $-\frac{3}{2}$)
\begin{center}
\begin{tikzpicture}[scale=0.85]
  \tkzTabInit[lgt=3,espcl=1.9]{$x$/1,$2x+3$/1,$x$/1,$x-2$/1,$\frac{x(x-2)}{2x+3}$/1.4}{$-\infty$,$-\frac{3}{2}$,$0$,$2$,$+\infty$}
  \tkzTabLine{,-,z,+,t,+,t,+,}
  \tkzTabLine{,-,t,-,z,+,t,+,}
  \tkzTabLine{,-,t,-,t,-,z,+,}
  \tkzTabLine{,-,d,+,z,-,z,+,}
\end{tikzpicture}
\end{center}
\end{expr}
\rem{Sous une valeur interdite, une double barre, pas un zéro. Chaque facteur $ax+b$ a le
signe de $a$ après sa valeur d'annulation.}
\end{exo}

\begin{exo}[Inéquations produit et quotient]
\begin{questions}[label=\alph*)]
  \item $9x^{2}-4=(3x-2)(3x+2)$ : on étudie $(3x+2)(3x-2)(x-2)\ge 0$.
\begin{center}
\begin{tikzpicture}[scale=0.85]
  \tkzTabInit[lgt=3.8,espcl=1.9]{$x$/1,$3x+2$/1,$3x-2$/1,$x-2$/1,$\left(9x^{2}-4\right)(x-2)$/1.2}{$-\infty$,$-\frac{2}{3}$,$\frac{2}{3}$,$2$,$+\infty$}
  \tkzTabLine{,-,z,+,t,+,t,+,}
  \tkzTabLine{,-,t,-,z,+,t,+,}
  \tkzTabLine{,-,t,-,t,-,z,+,}
  \tkzTabLine{,-,z,+,z,-,z,+,}
\end{tikzpicture}
\end{center}
        $S=\res{\left[-\frac{2}{3}\,;\frac{2}{3}\right]\cup\left[2\,;+\infty\right[}$
  \item $x^{2}-4=(x-2)(x+2)$ : valeurs interdites $-2$ et $2$.
\begin{center}
\begin{tikzpicture}[scale=0.85]
  \tkzTabInit[lgt=3,espcl=1.9]{$x$/1,$5-2x$/1,$x+2$/1,$x-2$/1,$\frac{5-2x}{x^{2}-4}$/1.4}{$-\infty$,$-2$,$2$,$\frac{5}{2}$,$+\infty$}
  \tkzTabLine{,+,t,+,t,+,z,-,}
  \tkzTabLine{,-,z,+,t,+,t,+,}
  \tkzTabLine{,-,t,-,z,+,t,+,}
  \tkzTabLine{,+,d,-,d,+,z,-,}
\end{tikzpicture}
\end{center}
        $S=\res{\left]-2\,;2\right[\cup\left[\frac{5}{2}\,;+\infty\right[}$
  \item $x^{2}-9=(x-3)(x+3)$ ; valeur interdite $-\frac{1}{3}$.
\begin{center}
\begin{tikzpicture}[scale=0.85]
  \tkzTabInit[lgt=3,espcl=1.9]{$x$/1,$x+3$/1,$3x+1$/1,$x-3$/1,$\frac{x^{2}-9}{3x+1}$/1.4}{$-\infty$,$-3$,$-\frac{1}{3}$,$3$,$+\infty$}
  \tkzTabLine{,-,z,+,t,+,t,+,}
  \tkzTabLine{,-,t,-,z,+,t,+,}
  \tkzTabLine{,-,t,-,t,-,z,+,}
  \tkzTabLine{,-,z,+,d,-,z,+,}
\end{tikzpicture}
\end{center}
        $S=\res{\left[-3\,;-\frac{1}{3}\right[\cup\left[3\,;+\infty\right[}$
  \item Valeur interdite $2$. On se ramène à un quotient comparé à $0$ :
        \begin{calculs}
        \dfrac{x+3}{2-x}\le 1 &\iff \dfrac{x+3-(2-x)}{2-x}\le 0\\[2pt]
                              &\iff \dfrac{2x+1}{2-x}\le 0
        \end{calculs}
\begin{center}
\begin{tikzpicture}[scale=0.85]
  \tkzTabInit[lgt=2.6,espcl=2]{$x$/1,$2x+1$/1,$2-x$/1,$\frac{2x+1}{2-x}$/1.4}{$-\infty$,$-\frac{1}{2}$,$2$,$+\infty$}
  \tkzTabLine{,-,z,+,t,+,}
  \tkzTabLine{,+,t,+,z,-,}
  \tkzTabLine{,-,z,+,d,-,}
\end{tikzpicture}
\end{center}
        $S=\res{\left]-\infty\,;-\frac{1}{2}\right]\cup\left]2\,;+\infty\right[}$
  \item Valeur interdite $-4$.
        \begin{calculs}
        \dfrac{2-x}{x+4}>3 &\iff \dfrac{2-x-3(x+4)}{x+4}>0\\[2pt]
                           &\iff \dfrac{-4x-10}{x+4}>0
        \end{calculs}
        $-4x-10$ s'annule en $-\frac{5}{2}$ ; son coefficient $-4$ est négatif.
\begin{center}
\begin{tikzpicture}[scale=0.85]
  \tkzTabInit[lgt=2.6,espcl=2]{$x$/1,$x+4$/1,$-4x-10$/1,$\frac{-4x-10}{x+4}$/1.4}{$-\infty$,$-4$,$-\frac{5}{2}$,$+\infty$}
  \tkzTabLine{,-,z,+,t,+,}
  \tkzTabLine{,+,t,+,z,-,}
  \tkzTabLine{,-,d,+,z,-,}
\end{tikzpicture}
\end{center}
        $S=\res{\left]-4\,;-\frac{5}{2}\right[}$
  \item Valeurs interdites $-2$ et $-3$. Numérateur :
        \begin{calculs}
        (x+1)(x+3)-(x-1)(x+2) &= \left(x^{2}+4x+3\right)-\left(x^{2}+x-2\right)\\
        (x+1)(x+3)-(x-1)(x+2) &= 3x+5
        \end{calculs}
        On étudie donc $\dfrac{3x+5}{(x+2)(x+3)}\le 0$.
\begin{center}
\begin{tikzpicture}[scale=0.85]
  \tkzTabInit[lgt=3,espcl=1.9]{$x$/1,$x+3$/1,$x+2$/1,$3x+5$/1,$\frac{3x+5}{(x+2)(x+3)}$/1.4}{$-\infty$,$-3$,$-2$,$-\frac{5}{3}$,$+\infty$}
  \tkzTabLine{,-,z,+,t,+,t,+,}
  \tkzTabLine{,-,t,-,z,+,t,+,}
  \tkzTabLine{,-,t,-,t,-,z,+,}
  \tkzTabLine{,-,d,+,d,-,z,+,}
\end{tikzpicture}
\end{center}
        $S=\res{\left]-\infty\,;-3\right[\cup\left]-2\,;-\frac{5}{3}\right]}$
\end{questions}
\rem{Une valeur interdite est toujours exclue de $S$, même avec $\le$ ou $\ge$. On ne
multiplie jamais par le dénominateur : son signe dépend de $x$.}
\end{exo}

\partie{Composition de polynômes}

\begin{exo}[Premières compositions]
\emph{Chaque $x$ de $P(x)$ est remplacé par $Q(x)$, entre parenthèses.}
\begin{questions}[label=\alph*)]
  \item $(P\circ Q)(x)=3(2x+1)+2=\res{6x+5}$
  \item $(P\circ Q)(x)=2(3x+2)+1=\res{6x+5}$
  \item \begin{calculs}
        (P\circ Q)(x) &= 3\left(2x^{2}+x-4\right)-2\\
        (P\circ Q)(x) &= \res{6x^{2}+3x-14}
        \end{calculs}
\end{questions}
\textbf{Comparaison.} En a) et b), on a échangé les rôles de $P$ et $Q$, et on obtient le
même polynôme. \res{C'est un hasard propre à ces deux fonctions} : en général
$P\circ Q\ne Q\circ P$ (exemple du cours, exercice 25).
\end{exo}

\begin{exo}[Composer et calculer]
\begin{questions}
  \item \begin{calculs}
        (P\circ Q)(x) &= (3x-1)^{2}+2(3x-1)-1\\
        (P\circ Q)(x) &= 9x^{2}-6x+1+6x-2-1\\
        (P\circ Q)(x) &= \res{9x^{2}-2}
        \end{calculs}
  \item \begin{calculs}
        (P\circ Q)(x) &= \left(x^{2}+1\right)^{2}-\left(x^{2}+1\right)\\
        (P\circ Q)(x) &= x^{4}+2x^{2}+1-x^{2}-1\\
        (P\circ Q)(x) &= \res{x^{4}+x^{2}}
        \end{calculs}
  \item \begin{questions}[label=\alph*)]
          \item \begin{calculs}
                (P\circ Q)(x) &= (x+4)^{2}-3\\
                (P\circ Q)(x) &= \res{x^{2}+8x+13}
                \end{calculs}
                donc $(P\circ Q)(-2)=4-16+13=\res{1}$.
          \item $Q(-2)=2$, puis $P(2)=4-3=\res{1}$ : on retrouve la même valeur.
        \end{questions}
\end{questions}
\rem{Pour une seule valeur, la seconde voie est souvent plus rapide.}
\end{exo}

\begin{exo}[Composer dans les deux sens]
\begin{questions}
  \item \begin{questions}[label=\alph*)]
          \item \begin{calculs}
                (P\circ Q)(x) &= \left(x^{2}-2x\right)^{2}+1\\
                (P\circ Q)(x) &= \res{x^{4}-4x^{3}+4x^{2}+1}\\[4pt]
                (Q\circ P)(x) &= \left(x^{2}+1\right)^{2}-2\left(x^{2}+1\right)\\
                (Q\circ P)(x) &= x^{4}+2x^{2}+1-2x^{2}-2\\
                (Q\circ P)(x) &= \res{x^{4}-1}
                \end{calculs}
          \item $P$ et $Q$ sont de degré \res{2}, $P\circ Q$ est de degré \res{4}.
        \end{questions}
  \item \begin{calculs}
        (P\circ Q)(x) &= 2\left(\sqrt{2}\,x+1\right)^{2}-\left(\sqrt{2}\,x+1\right)\\
        (P\circ Q)(x) &= 2\left(2x^{2}+2\sqrt{2}\,x+1\right)-\sqrt{2}\,x-1
          &&\justif{$\left(\sqrt{2}\,x\right)^{2}=2x^{2}$}\\
        (P\circ Q)(x) &= 4x^{2}+4\sqrt{2}\,x+2-\sqrt{2}\,x-1\\
        (P\circ Q)(x) &= \res{4x^{2}+3\sqrt{2}\,x+1}
        \end{calculs}
\end{questions}
\rem{Les deux composées du 1\textdegree) sont différentes : l'ordre compte.}
\end{exo}

\end{document}
